% GENERATED FILE. Edit canonical lessons, then run npm run generate:books.
% canonical-source-hash: fnv1a64:fb294829aa56f5fc
% canonical-lessons: TE-C276-rasayanasastram, TE-C276-jivasastram, TE-C276-arthasastram, TE-C276-vyakaranam, TE-C276-nallaballa

\chapter{Chemistry, Biology, Economics, Grammar, Blackboard}
\label{ch:te-chemistry-biology-economics-grammar-blackboard}
\hlchaptermodality{276}

\begin{chapteropening}
\textbf{By the end of this chapter:} \emph{I can say chemistry, biology, economics, grammar and a blackboard.}

Complete the last lesson of chapter 276: I can say chemistry, biology, economics, grammar and a blackboard.
\end{chapteropening}

\section[\texorpdfstring{\te{రసాయనశాస్త్రం} (rasāyanaśāstraṁ)}{రసాయనశాస్త్రం (rasāyanaśāstraṁ)}]{\te{రసాయనశాస్త్రం} (rasāyanaśāstraṁ) — chemistry}

\label{lesson:TE-C276-rasayanasastram}



\begin{quote}
Before the new one: say the Telugu for music, then the Telugu for a play.
\end{quote}

\subsection*{You'll want to know: \te{రసాయనశాస్త్రం}}
\textbf{\te{రసాయనశాస్త్రం}} — \emph{rasāyanaśāstraṁ} — \textquotedblleft{}chemistry\textquotedblright{}.

\begin{grammarlens}[title={What you've built}]
One more word for this chapter.
\end{grammarlens}

\subsection*{Your turn}
\begin{itemize}
\raggedright
  \item \emph{Say it:} \emph{rasāyanaśāstraṁ}
  \item \emph{Say it:} \emph{rasāyanaśāstraṁ}, once more
  \item \emph{Say it:} say \emph{rasāyanaśāstraṁ}
  \item \emph{Recall:} say the Telugu for music, then the Telugu for a play, then say \emph{rasāyanaśāstraṁ} again
\end{itemize}

\begin{tcolorbox}[breakable,colback=teal!4,colframe=teal!35!black,title={Before you move on}]
What does \te{రసాయనశాస్త్రం} mean? (\textquotedblleft{}Chemistry\textquotedblright{}.) Say it once more.
\end{tcolorbox}
\section[\texorpdfstring{\te{జీవశాస్త్రం} (jīvaśāstraṁ)}{జీవశాస్త్రం (jīvaśāstraṁ)}]{\te{జీవశాస్త్రం} (jīvaśāstraṁ) — biology}

\label{lesson:TE-C276-jivasastram}



\begin{quote}
Before the new one: say the Telugu for a play, then the Telugu for chemistry.
\end{quote}

\subsection*{You'll want to know: \te{జీవశాస్త్రం}}
\textbf{\te{జీవశాస్త్రం}} — \emph{jīvaśāstraṁ} — \textquotedblleft{}biology\textquotedblright{}.

\begin{grammarlens}[title={What you've built}]
One more word for this chapter.
\end{grammarlens}

\subsection*{Your turn}
\begin{itemize}
\raggedright
  \item \emph{Say it:} \emph{jīvaśāstraṁ}
  \item \emph{Say it:} \emph{jīvaśāstraṁ}, once more
  \item \emph{Say it:} say \emph{jīvaśāstraṁ}
  \item \emph{Recall:} say the Telugu for a play, then the Telugu for chemistry, then say \emph{jīvaśāstraṁ} again
\end{itemize}

\begin{tcolorbox}[breakable,colback=teal!4,colframe=teal!35!black,title={Before you move on}]
What does \te{జీవశాస్త్రం} mean? (\textquotedblleft{}Biology\textquotedblright{}.) Say it once more.
\end{tcolorbox}
\section[\texorpdfstring{\te{అర్థశాస్త్రం} (arthaśāstraṁ)}{అర్థశాస్త్రం (arthaśāstraṁ)}]{\te{అర్థశాస్త్రం} (arthaśāstraṁ) — economics}

\label{lesson:TE-C276-arthasastram}



\begin{quote}
Before the new one: say the Telugu for chemistry, then the Telugu for biology.
\end{quote}

\subsection*{You'll want to know: \te{అర్థశాస్త్రం}}
\textbf{\te{అర్థశాస్త్రం}} — \emph{arthaśāstraṁ} — \textquotedblleft{}economics\textquotedblright{}.

\begin{grammarlens}[title={What you've built}]
One more word for this chapter.
\end{grammarlens}

\subsection*{Your turn}
\begin{itemize}
\raggedright
  \item \emph{Say it:} \emph{arthaśāstraṁ}
  \item \emph{Say it:} \emph{arthaśāstraṁ}, once more
  \item \emph{Say it:} say \emph{arthaśāstraṁ}
  \item \emph{Recall:} say the Telugu for chemistry, then the Telugu for biology, then say \emph{arthaśāstraṁ} again
\end{itemize}

\begin{tcolorbox}[breakable,colback=teal!4,colframe=teal!35!black,title={Before you move on}]
What does \te{అర్థశాస్త్రం} mean? (\textquotedblleft{}Economics\textquotedblright{}.) Say it once more.
\end{tcolorbox}
\section[\texorpdfstring{\te{వ్యాకరణం} (vyākaraṇaṁ)}{వ్యాకరణం (vyākaraṇaṁ)}]{\te{వ్యాకరణం} (vyākaraṇaṁ) — grammar}

\label{lesson:TE-C276-vyakaranam}



\begin{quote}
Before the new one: say the Telugu for biology, then the Telugu for economics.
\end{quote}

\subsection*{You'll want to know: \te{వ్యాకరణం}}
\textbf{\te{వ్యాకరణం}} — \emph{vyākaraṇaṁ} — \textquotedblleft{}grammar\textquotedblright{}.

\begin{grammarlens}[title={What you've built}]
One more word for this chapter.
\end{grammarlens}

\subsection*{Your turn}
\begin{itemize}
\raggedright
  \item \emph{Say it:} \emph{vyākaraṇaṁ}
  \item \emph{Say it:} \emph{vyākaraṇaṁ}, once more
  \item \emph{Say it:} say \emph{vyākaraṇaṁ}
  \item \emph{Recall:} say the Telugu for biology, then the Telugu for economics, then say \emph{vyākaraṇaṁ} again
\end{itemize}

\begin{tcolorbox}[breakable,colback=teal!4,colframe=teal!35!black,title={Before you move on}]
What does \te{వ్యాకరణం} mean? (\textquotedblleft{}Grammar\textquotedblright{}.) Say it once more.
\end{tcolorbox}
\section[\texorpdfstring{\te{నల్లబల్ల} (nallaballa)}{నల్లబల్ల (nallaballa)}]{\te{నల్లబల్ల} (nallaballa) — a blackboard}

\label{lesson:TE-C276-nallaballa}



\begin{quote}
Before the new one: say the Telugu for economics, then the Telugu for grammar.
\end{quote}

\subsection*{You'll want to know: \te{నల్లబల్ల}}
\textbf{\te{నల్లబల్ల}} — \emph{nallaballa} — \textquotedblleft{}a blackboard\textquotedblright{}.

\begin{grammarlens}[title={What you've built}]
One more word for this chapter.
\end{grammarlens}

\subsection*{Your turn}
\begin{itemize}
\raggedright
  \item \emph{Say it:} \emph{nallaballa}
  \item \emph{Say it:} \emph{nallaballa}, once more
  \item \emph{Say it:} say \emph{nallaballa}
  \item \emph{Recall:} say the Telugu for economics, then the Telugu for grammar, then say \emph{nallaballa} again
\end{itemize}

\begin{tcolorbox}[breakable,colback=teal!4,colframe=teal!35!black,title={Before you move on}]
What does \te{నల్లబల్ల} mean? (\textquotedblleft{}A blackboard\textquotedblright{}.) Say it once more.
\end{tcolorbox}
